\documentclass[../../main.tex]{subfiles}
\begin{document}

{\bf [解]}
\paragraph{(1)}

$a=\sqrt{2}$とする．
$1 < \sqrt{2} < 2$ から，$\langle a \rangle = \sqrt{2}-1$ である．
従って
\begin{align}
  a_1 = \sqrt{2}-1, a_2 = \langle \dfrac{1}{\sqrt{2}-1} \rangle = \sqrt{2}-1
\end{align}
だから以下帰納的に$a_n = \sqrt{2}-1$である．

\paragraph{(2)}

まず $a_1 = \langle a \rangle$ だから，これが$a$となるために
\begin{align}
  \dfrac{1}{3} \le a < 1 \label{eq:1}
\end{align}
が必要である．このもとで漸化式から
\begin{align}
  a_2 = \left\langle \dfrac{1}{a_1} \right\rangle = \left\langle \dfrac{1}{a} \right\rangle = a \label{eq:2}
\end{align}
だから$\dfrac{1}{a}$以下の最大の整数を$m \in \mathbb{Z}_{\ge 0}$ として
\begin{align}
  \dfrac{1}{a} - m = a \quad \therefore a^2 + ma - 1 = 0
\end{align}
である．この方程式を解いて
\begin{align}
  a = \dfrac{-m \pm \sqrt{m^2+4}}{2} \label{eq:3}
\end{align}
だから，\cref{eq:1}に代入して
\begin{align}
  & \dfrac{2}{3} \le -m \pm \sqrt{m^2+4} < 2 \\
  \therefore
  & m + \dfrac{2}{3} \le \pm \sqrt{m^2+4} < m + 2
\end{align}
である．
$m \ge 0$ から複合正が正しく\cref{eq:3}は
\begin{align}
  a = \dfrac{-m + \sqrt{m^2+4}}{2} \label{eq:4}
\end{align}
である．

この元で不等式の各項0以上だから2乗してよく，
\begin{align}
  \dfrac{4}{9} m + \dfrac{4}{9} \le 4 < 4m + 4
  \quad \therefore
  0 < m \le \dfrac{8}{3}
\end{align}
を得る．
$m \in \mathbb{Z}_{\ge 0}$ よりこれを満たすのは $m = 1, 2$ である．
\cref{eq:4}に代入して
\begin{align}
  a = \dfrac{-1+\sqrt{5}}{2}, -1+\sqrt{2}
\end{align}
である．この時帰納的に$a_n=a$となるから求める$a$はこの二つである．

\paragraph{(3)}

$p$ と $q$ は互いに素であるとする．
$a_n \neq 0$ かつ $a_{n+1}=0$ の時条件 (ii)から $\dfrac{1}{a_n} \in \mathbb{Z}$ である．

$q=1$ の時は明らかに題意は成立するから以下 $q \ge 2$ とする．
$p$ を $q$ で割った商を $j_1$，あまりを $r_1$ として
\begin{align}
  p = j_1 q + r_1 \quad (0 \le r_1 < q) \label{eq:5}
\end{align}
とかく．

$r_1=0$ の時は $a_1 = 0$ となり題意は成立するから以下 $r_1 \neq 0$ とする．
この時
\begin{align}
  \left\langle a \right\rangle
  = \left\langle \dfrac{p}{q} \right\rangle
  = \left\langle j_1 + \dfrac{r_1}{q} \right\rangle
  = \dfrac{r_1}{q} \quad (\neq 0)
\end{align}
となる．$q$ を $r_1$ で割った商とあまりを，$j_2, r_2$ とすると
\begin{align}
  q
  &= j_2 r_1 + r_2 \quad (0 \le r_2 < r_1) \\
  \therefore
  a_2
  &= \left\langle \dfrac{1}{a_1} \right\rangle \\
  &= \left\langle \dfrac{q}{r_1} \right\rangle \\
  &= \dfrac{r_2}{r_1}
\end{align}
$r_2=0$ ならば $a_2=0$，$r_2 \neq 0$ ならば $r_1$ を $r_2$ で割った余りを $r_3$ とする．以下同様に，$m \in \mathbb{N}$ に対し，
$r_m = 0$ なら $a_m = 0$，$r_m \neq 0$ なら $r_{m-1}$ を $r_m$ で割ったあまりを $r_{m+1}$ とする．（ただし $r_0 = q$）
ここで，$q$ 以下の自然数 $l$ で，$r_l = 0$ となるものが存在しないと仮定すると，
\begin{align}
  q > r_1 > r_2 > \cdots > r_q \ge 1
\end{align}
なる自然数 $r_1, r_2, \dots, r_q$ が得られて矛盾する．
よって $r_l = 0$ となるものが存在しこの時 $q$ 以上の全ての自然数に対し，$a_n = 0$ となる．
以上から題意は示された．

\end{document}
